\section{Reproducible experimental methodology}
\subsection{Study design and staged feasibility gates}
The proposed research is a prospective, controlled comparative study. Gate~0 establishes whether the actual coating emits measurable light at the proposed excitation wavelength, and whether the UV enclosure is safe. Gate~1 establishes optical repeatability and ground-truth reproducibility. Gate~2 develops three algorithmic approaches with separate training and validation boards. Gate~3 freezes all thresholds and models, then measures performance on completely held-out boards and ideally held-out production lots. Gate~4 evaluates production robustness, human override and timing over a sufficiently long representative sequence. Every stage has an explicit failure condition: if critical defects are not optically observable, further neural training alone is not a satisfactory resolution.

\subsection{Specimen and lot sampling}
Each physical PCB must have a unique immutable \texttt{specimen\_id}; each production/process batch has \texttt{lot\_id}. Every image receives an \texttt{image\_id} linking it to coating material, preparation, light arrangement, camera parameters, board orientation, cure status, capture time and annotation revision. Acquire multiple images from the same physical board to test optical repeatability but keep all images from that board \emph{within one data partition}. Where enough lots exist, use grouped splits by lot (train lots, validation lots, independent test lots). If only one lot exists, report that generalization across lots has not yet been tested. The companion \texttt{dataset\_register.csv} provides required tracking columns.

An initial \emph{planning target} could be approximately 100--200 physical boards across at least three visually distinct collection sessions or lots; this is not a power calculation and should not be presented as automatically sufficient for a rare-defect detection claim. Exact physical board counts must follow expected defect prevalence, the minimum acceptable sensitivity, the feasible number of independently confirmed defect boards, and confidence-interval targets. If natural faults are rare, create deliberately defective test boards with controlled omission, overspray, bubble-generating process variation or masking failures, subject to manufacturing and safety approval. Artificial examples should supplement, not replace, naturally occurring failures in the final production validation. A synthetic pixel hole applied in software is a useful stress test but not a true coating defect specimen.

\subsection{Defect taxonomy and acceptance thresholds}
\begin{table}[htbp]\centering\footnotesize
\caption{Candidate annotation and decision ontology. Inspection feasibility and acceptance limits are to be established, not assumed.}\label{tab:ontology}
\begin{tabularx}{\textwidth}{p{30mm}p{25mm}X}\toprule
Class & Label representation & Independent confirmation / concern\\\midrule
Missing coating & Pixel polygon or region instance & UV evidence plus expert/alternate inspection; distinguish intentional mask-out\\
Excess coating / overspray & Pixel region crossing forbidden zone & Coating plan revision and expert confirmation\\
Bubble / void & Instance region, size and location & Microscopy or expert adjudication; fluorescent dark spots are not all bubbles\\
Dewetting & Boundary polygon, area and context & Process engineer confirmation; often resembles absence or capillary effect\\
Capillary pooling & Region plus component proximity & Fluorescence intensity alone may not quantify thickness\\
Crack / scratch & Instance or line segment & Alternate-angle white/UV imaging; distinguish surface artefacts\\
Occluded / not observable & Exclusion mask & Must be handled as unknown, not labelled sound coating\\
\bottomrule\end{tabularx}
\end{table}

A defect is \emph{actionable} only if a documented customer, internal drawing or authorized acceptance specification makes it rejectable under the applicable product class. Record minimum size, permitted locations, any exceptions for component leads and acceptable bubble conditions. The goal is traceable compliance with a validated inspection plan, not the maximum number of bright/dark pixels detected. The full IPC-A-610J and coating plan must be consulted by a suitably qualified person before production rules are finalized \citep{ipca610j,ipc2024release}.

\subsection{Optical acquisition protocol}
For each run, record LED electrical settings and ideally calibrated irradiance (or at least an inter-session stable radiometric monitor), exposure time, analog gain, white balance state, image bit depth, lens aperture/focus, working distance, board-to-camera distance and filter identification. Use lossless raw or minimally processed image data wherever practical; JPEG compression can distort tiny bubble contours. Capture a dark frame and fluorescence reference after equipment warm-up. Image a bare uncoated board or verified keep-out surface to estimate background auto-fluorescence. Confirm that the visible-emission band is dominant after optical filtering \citep{crowther2022filters}.

The nominal factorial optical study should vary lighting configuration (top-only, two opposing sides, four-source combined, sequential directional exposures), exposure (at least three levels avoiding saturation), and image scale (for example ~5, ~8, ~12 MP optical equivalents or calibrated downsampled conditions). Do not use a downsampled 12~MP image as complete substitute for testing an actual 5~MP sensor: lens pixel size, color filter array, noise and readout differ. Each configuration should include multiple independent boards in randomized capture order; sources of non-independence should be modelled or accounted for when estimating uncertainty. Full-factorial designs may be expensive; staged screening followed by a focused confirmatory design is acceptable if predeclared.

\subsection{Ground truth and observer agreement}
Ground truth must be independent of the algorithms being evaluated. Two trained inspectors should independently label a predefined subset, blinded to model output, then reconcile disagreement using qualified engineering or microscopy review. Store \emph{original labels}, adjudicated labels, defect class, criticality, size and ``uncertain'' states. Record inter-annotator agreement on actionable defect identification and segmentation where meaningful; IoU/Dice between masks and agreement on board disposition measure different things. Some ground-truth evidence may be obtained by selective re-inspection at higher magnification, microscopy or OCT. If two experts cannot reliably agree on an image, model failure against one arbitrary mask should not be misinterpreted as a purely algorithmic problem.

\subsection{Classical baseline A}
Develop a fully documented corrected-image threshold pipeline using Otsu/adaptive thresholds as alternatives, threshold channel selected using validation data, fixed morphology based on millimetre-scale rules, and connected-component grouping \citep{otsu1979}. Register the board geometry using fiducial-driven affine/projective fitting and fail closed on poor registration \citep{fischler1981ransac}. Evaluate both a \emph{simple threshold-only} system and a \emph{classical plus reference masks} system. Without this latter baseline, improved performance of the hybrid could merely come from adding the reference plan rather than introducing AI.

\subsection{Learned baseline B}
Start with a reproducible U-Net variant and, if needed, SegFormer-B0 or MobileNetV3-based segmentation \citep{ronneberger2015unet,xie2021segformer,howard2019mobile}. Record fixed random seeds, model input resolution, weight initialization, class balance, augmentation, loss definition, optimizer, learning rate and early stopping. Set architecture and input resolution before viewing held-out test data. Training data may be augmented for lighting shifts, rotations and focus/noise within realistic ranges; prohibit spatial warps that invalidate defect sizes or move keep-out masks independently of the image. Use the same physical-board split and comparable augmentation conditions for competing models.

\subsection{Hybrid candidate C and ablations}
Run exactly the same optical acquisitions through registration, a learned coating probability map and geometry-aware required/forbidden region analysis as proposed in Section~6. Predeclare rules for a thresholded missing-coating area, largest missing island, proximity to critical geometry and uncertain-region abstention. Experimental factors are the presence or absence of correction, registration, reference map, learned branch and alternate illumination. A useful factorial ablation table includes: classical only, classical+reference, CNN only, CNN+reference, CNN+reference+correction, and model with multiple exposures. Report where improved results arise from a particular stage. A learned feature extractor could be replaced with PatchCore or EfficientAD in a secondary comparison, but avoid an unmanageable unprioritized model zoo \citep{roth2022patchcore,batzner2024efficientad}.

\subsection{Defect localization and fair matching}
Set a physical detection threshold. For a predicted defect instance $P_i$ and independently annotated ground-truth instance $G_j$, match using intersection-over-union (IoU) threshold, centroid tolerance or any-domain-specific criterion chosen \emph{a priori}. Tiny line defects may need a tolerance in millimetres rather than IoU alone. Use one-to-one matching between predictions and ground truth to avoid double counting. Summarize counts by defect class, size, location, illumination and criticality. Count a missed actionable physical defect as a false negative even if the image classifier correctly identifies another defect elsewhere on the same board.

\subsection{Production timing and system verification}
The complete cycle includes trigger-to-exposure, optional multiple exposures, camera transfer, acquisition validation, geometric registration, any inference tiling, defect aggregation, logging/database output, PLC result handshake and safety/fault reporting. Measure median, 95th and 99th percentile times across sustained runs, along with cold start, model reload, CPU load, RAM, thermal throttling and recovery after camera disconnection. A nominal inference-only speed on a GPU-equipped laptop cannot be copied to a used Intel office PC. The workstation's actual takt requirement must be supplied by production engineering and the system benchmarked against it \citep{onnxproviders,openvino2026}.

\subsection{Reproducibility and leakage checks}
Before each experiment: (1) confirm dataset hash and physical-board split; (2) confirm all label revisions; (3) record software Git commit, Python/CV/ML library versions and hardware; (4) freeze train/validation/test membership; (5) save calibrated acquisition parameters; (6) archive complete confusion matrices and false-decision images. Training examples acquired from the same specimen as the test image are an especially likely source of inflated accuracy. Avoid using a test specimen as the reference template for its own detection or including its other exposures in training.
